%%
%% CSD Assignment 1 - Final Technical Report Template
%% Based on the ACM "acmart" manuscript (single-column) style.
%% Fill in every section marked [FILL IN]. Delete the guidance
%% text (in italics) once you have written your own content.
%%
\documentclass[manuscript,screen,review]{acmart}

\AtBeginDocument{%
  \providecommand\BibTeX{{Bib\TeX}}}

%% ------------------------------------------------------------------
%% This is a COURSE ASSIGNMENT report, not an ACM publication, so the
%% rights-management / DOI / ISBN block from the original ACM sample
%% is not required. It is left commented out below in case your
%% instructor asks you to keep the ACM-style header block.
%% ------------------------------------------------------------------
% \setcopyright{acmlicensed}
% \acmConference[CSD Assignment 1]{Course: Computer Science and Design}{2026}{Your Institute}
% \acmISBN{}

\begin{document}

%% =================== TITLE & AUTHOR BLOCK =========================
\title{AI-based Village Pond Planning System}
\subtitle{CSD Assignment 1 -- Final Technical Report}

\author{Your Full Name}
\affiliation{%
  \institution{Your Institute / Department}
  \city{Your City}
  \country{India}}
\email{your.email@institute.ac.in}

%% Add more \author{} + \affiliation{} + \email{} blocks here if this
%% is a group project. Keep one block per team member.

\renewcommand{\shortauthors}{Your Surname et al.}

%% =================== ABSTRACT ======================================
\begin{abstract}
\textit{[FILL IN -- 150--250 words.]} Summarize the problem (rural
water conservation and pond-site selection), the approach (geospatial
+ rainfall data analysis exposed through a web application), the key
system components you built, and the main results (e.g., number of
candidate sites evaluated, accuracy/usefulness of the recommended
pond depth and storage capacity, and any performance figures for the
system itself).
\end{abstract}

\keywords{Village Pond Planning, Geospatial Analysis, Rainwater
  Harvesting, Catchment Estimation, Web GIS, REST APIs}

\maketitle

%% =====================================================================
\section{Introduction}
\label{sec:intro}
\textit{[MUST BE INCLUDED]} Introduce the problem of rural water conservation
and why pond-based rainwater harvesting is an effective solution.
State the goal of the assignment (an AI/geospatial-assisted web
application that recommends suitable pond locations) and give a short
roadmap of the rest of the report.

\subsection{Motivation}
\textit{[FILL IN]} Explain, in your own words, why manual site
selection is hard (terrain complexity, lack of centralized rainfall
data, unavailability of land-ownership records, etc.) and why an
automated tool helps village administrators.

\subsection{Scope of the Project}
\textit{[MUST BE INCLUDED]} List, in one paragraph or a short bullet list, what
your system does and does \emph{not} cover (e.g., "estimates catchment
area and expected runoff for a user-selected point; does not perform
structural/civil engineering design of the pond embankment").

%% =====================================================================
\section{Problem Statement and Requirements}
\label{sec:requirements}
\textit{[MUST BE INCLUDED]} Restate the assignment's functional requirements in
your own words and map each to the module of your system that
satisfies it. A table works well here.

\begin{table}[h]
  \caption{Functional requirements and where they are implemented}
  \label{tab:requirements}
  \begin{tabular}{p{3.6cm}p{3.6cm}}
    \toprule
    Requirement & Implemented in \\
    \midrule
    Satellite imagery display & \textit{[module/file]} \\
    Contour map visualization & \textit{[module/file]} \\
    Available-land identification & \textit{[module/file]} \\
    Catchment area estimation & \textit{[module/file]} \\
    Historical rainfall query & \textit{[module/file]} \\
    Runoff volume estimation & \textit{[module/file]} \\
    Pond depth / storage recommendation & \textit{[module/file]} \\
    Combined overlay / results view & \textit{[module/file]} \\
    \bottomrule
  \end{tabular}
\end{table}

\subsection{Non-Functional Requirements}
\textit{[MUST BE INCLUDED]} State performance, scalability, availability,
usability, and security requirements you set for yourself, e.g.
response time for a query, number of concurrent users the system
should tolerate, expected uptime, browser support, etc.


%% =====================================================================
\section{System Architecture and High-Level Design}
\label{sec:architecture}
\textit{[MUST BE INCLUDED]} Describe the overall architecture: client, server,
database, and external API layers. Include an architecture diagram
(Figure~\ref{fig:architecture}) showing how data flows from the map
UI to the backend, to the elevation/rainfall APIs, to the database,
and back.

\begin{figure}[h]
  \centering
  % \includegraphics[width=\linewidth]{architecture-diagram}
  \caption{High-level system architecture.}
  \label{fig:architecture}
  \Description{Replace with an accurate description of your
    architecture diagram for accessibility.}
\end{figure}

\subsection{Technology Stack}
\textit{[MUST BE INCLUDED]} List the actual stack you used (language, web
framework, database, frontend library/framework, mapping library,
and which elevation/rainfall APIs you queried). Justify any deviation
from the suggested stack in the assignment brief.

%% =====================================================================
\section{Methodology}
\label{sec:methodology}
[MUST BE INCLUDED] This section should explain the actual algorithms and data pipelines
you implemented -- this is the core of the "Terrain and Catchment
Analysis" evaluation criterion, so be specific and quantitative.

\subsection{Terrain and Elevation Analysis}
\textit{[FILL IN]} Explain how you obtained and processed the
Digital Elevation Model (DEM) / elevation API data, how contour lines
were generated, and how you used OpenCV (or your chosen library) for
any image/terrain processing.

\subsection{Catchment Area Delineation}
\textit{[FILL IN]} Explain the algorithm used to estimate the
catchment area contributing runoff to a candidate point (e.g.,
flow-direction/flow-accumulation from the DEM, or a simplified
buffer/slope-based heuristic). State assumptions and limitations.

\subsection{Rainfall Data Integration}
\textit{[FILL IN]} Describe which rainfall API(s) you queried (IMD,
Open-Meteo, NASA POWER, etc.), the time range of historical data
used, and how you aggregated it into annual/seasonal statistics.

%% =====================================================================
\section{Implementation}
\label{sec:implementation}

\subsection{Backend and API Design}
\textit{[MUST BE INCLUDED]} List your REST/GraphQL endpoints (a table of
\texttt{METHOD /path -- purpose} works well), authentication approach
if any, and how errors from third-party APIs are handled.

\subsection{Frontend and Visualization}
\textit{[MUST BE INCLUDED]} Describe the map interface, the visualization
library used, and how the results overlay (site, catchment boundary,
rainfall stats, runoff volume, pond dimensions) is presented to the
user. Include a screenshot.

\begin{figure}[h]
  \centering
  % \includegraphics[width=\linewidth]{ui-screenshot}
  \caption{Application interface showing the results overlay.}
  \label{fig:ui}
  \Description{Replace with an accurate description of the screenshot.}
\end{figure}

\subsection{Database and Storage}
\textit{[FILL IN]} Explain how spatial data, rainfall records, and
computed recommendations are persisted and queried.

%% =====================================================================
\section{CSD Themes and Topics Applied in the Project}
\label{sec:csd-themes}
\textit{[MUST BE INCLUDED] This section is where you explicitly connect your
implementation back to core Computer System Design (CSD)
concepts taught in the course. For every theme you list, name the
concrete component of your system where it appears and justify why
you made that design choice (not just that the feature exists).
Add or remove rows/subsections to match what you actually built --
do not claim a theme you did not implement.}

\begin{table*}[t]
  \caption{Mapping of core CSD themes/topics to their use in this project}
  \label{tab:csd-themes}
  \begin{tabular}{p{2.8cm}p{5.2cm}p{6.5cm}}
    \toprule
    CSD Theme/Topic & Where used in the project & Justification / design rationale \\
    \midrule
    API Design (REST) & \textit{[e.g., \texttt{/api/catchment}, \texttt{/api/rainfall}]} & \textit{[Why REST, versioning, request/response contract]} \\
    Load Balancing & \textit{[e.g., Nginx in front of N app instances]} & \textit{[Why needed -- concurrent map queries, API rate limits, etc.]} \\
    Caching & \textit{[e.g., cached rainfall API responses / tile cache]} & \textit{[Reduces external API calls, improves latency]} \\
    Database Indexing / Query Optimization & \textit{[e.g., spatial index on site coordinates]} & \textit{[Speeds up geospatial queries]} \\
    Concurrency / Asynchronous Processing & \textit{[e.g., async calls to elevation + rainfall APIs]} & \textit{[Avoids blocking the UI while multiple APIs are queried]} \\
    Microservices vs. Monolith & \textit{[state which you chose]} & \textit{[Trade-offs you considered]} \\
    Design Patterns & \textit{[e.g., Factory for API clients, Strategy for runoff models]} & \textit{[Why the pattern fit the problem]} \\
    Authentication / Authorization & \textit{[if applicable]} & \textit{[Who can trigger analyses / view admin data]} \\
    Containerization / Deployment & \textit{[e.g., Docker, docker-compose]} & \textit{[Reproducible environment, ease of grading/demo]} \\
    Error Handling and Resilience & \textit{[e.g., retries/fallbacks when a rainfall API is down]} & \textit{[System should degrade gracefully]} \\
    Algorithms and Complexity & \textit{[e.g., flow-accumulation algorithm complexity]} & \textit{[Why this algorithm was chosen over alternatives]} \\
    Testing Strategy & \textit{[unit/integration/system tests -- see Section~\ref{sec:testing}]} & \textit{[What confidence this gives in correctness]} \\
    Version Control / CI-CD & \textit{[e.g., GitHub Actions pipeline]} & \textit{[How it kept the team in sync / caught regressions]} \\
    \bottomrule
  \end{tabular}
\end{table*}

\textit{[FILL IN]} Add a short paragraph after the table walking
through the two or three themes you consider most significant to
this project and explaining, in more depth, the specific engineering
decision behind each (e.g., "We added a load balancer because X" or
"We chose not to implement microservices because Y").

%% =====================================================================
\section{Results and Evaluation}
\label{sec:results}
\textit{[MUST BE INCLUDED]} Present representative results: example
village/site with its computed catchment area, rainfall statistics,
estimated runoff, and recommended pond dimensions. Include maps or
charts (Figure~\ref{fig:results}) and, if possible, compare your
estimated values against any ground-truth or published reference
figures.

\begin{figure}[h]
  \centering
  % \includegraphics[width=\linewidth]{results-overlay}
  \caption{Example results overlay for a selected village.}
  \label{fig:results}
  \Description{Replace with an accurate description of the results
    figure.}
\end{figure}

\subsection{Performance}
\textit{[MUST BE INCLUDED]} Report response times for key operations (e.g.,
time to fetch rainfall data, time to compute catchment area for a
given site) and how the system behaves under load if you tested that.

%% =====================================================================
\section{Discussion and Limitations}
\label{sec:discussion}
\textit{[FILL IN]} Reflect honestly on what worked well, what did
not, data-quality issues (e.g., coarse DEM resolution, gaps in
rainfall records), and any simplifying assumptions that affect the
accuracy of the recommendations.

%% =====================================================================
\section{AI Tool Usage Declaration}
\label{sec:ai-usage}
\textit{[MUST BE INCLUDED -- required by the assignment's LLM Usage Policy.]}
State which AI tools (ChatGPT, Claude, Gemini, GitHub Copilot, etc.)
were used, for which specific tasks (brainstorming, debugging,
documentation, refactoring, etc.) like AI was used to reformat this tempalte, and confirm that all AI-assisted
code was reviewed, understood, and modified by the team before
submission. Do not leave this section blank.


%% =====================================================================
\appendix

\section{Source Code and Repository}
\textit{[MUST BE INCLUDED]} Provide the repository link and a short note on
its folder structure.

\end{document}
\endinput
